% TOP_PROBLEM: {"id": 9600003, "title": "Exchangeability in the mean-zero exclusion process", "category_id": 19, "category": {"id": 19, "name": "probability", "display_name": "Probability", "description": "Problems involving probability theory, stochastic processes, and random structures.", "slug": "probability", "order_index": 19, "created_at": "2026-07-31T15:26:25.671Z"}, "status": "partially_solved", "problem_number": "AMR-095-0003", "set_id": 13, "statusQualification": "Restores the source’s introductory irreducibility assumption and explicitly interprets a finite zero mean by absolute first-moment integrability."}
% FIRST_POSTED: 2026-10-01
% ENTRY_KIND: statement_only
% UnsolvedMath ID 9600003; source catalogue code AMR-095-0003
% !TeX program = lualatex
% Definition and sources only; this file is not an LLM solution attempt.
\documentclass[11pt,a4paper]{article}
\usepackage[margin=25mm]{geometry}
\usepackage{fontspec}
\setmainfont{lmroman10-regular.otf}[BoldFont=lmroman10-bold.otf,
 ItalicFont=lmroman10-italic.otf,BoldItalicFont=lmroman10-bolditalic.otf]
\usepackage{amsmath,amssymb,mathtools}
\usepackage{unicode-math}
\setmathfont{latinmodern-math.otf}
\AtBeginDocument{\let\setminus\smallsetminus}
\usepackage{xurl,hyperref}
\hypersetup{colorlinks=true,urlcolor=blue}
\setcounter{secnumdepth}{0}
\setlength{\parindent}{0pt}
\setlength{\parskip}{5pt}
\setlength{\emergencystretch}{3em}
\begin{document}
\section*{Exchangeability in the mean-zero exclusion process}\label{9600003}
{\small UnsolvedMath ID 9600003 · AMR-095-0003 · Probability · AMR Open Problem Lists}
\subsection{Definitions and mathematical statement}
Fix $d\ge1$ and a probability kernel $p:\mathbb Z^d\to[0,1]$ with finite first absolute moment and zero mean:
\[
\sum_z p(z)=1,\qquad \sum_z\|z\|p(z)<\infty,
\qquad \sum_z z p(z)=0.
\]
Assume that the random walk with jumps distributed as $p$ is irreducible on $\mathbb Z^d$. In the exclusion process, each ordered pair $(x,x+z)$ has a rate-$p(z)$ clock. A ring moves a particle from $x$ to $x+z$ exactly when the former site is occupied and the latter is empty. At most one particle occupies a site.

A probability measure $\mu$ on $\{0,1\}^{\mathbb Z^d}$ is stationary if it is preserved by this evolution. It is exchangeable if its law is unchanged by every permutation of lattice sites which fixes all but finitely many sites. Must every stationary measure be exchangeable?

Equivalently, must there exist a probability measure $\nu$ on $[0,1]$ such that
\[
\mu=\int_0^1\bigotimes_{x\in\mathbb Z^d}\operatorname{Bernoulli}(\rho)\,\nu(d\rho)?
\]
Thus the proposed classification allows a random global density but requires conditional independence and identical occupation probabilities once that density is chosen. No spatial invariance hypothesis is placed on $\mu$: it is part of the proposed conclusion. The zero-mean condition permits asymmetric jump laws and does not require $p(z)=p(-z)$. Irreducibility prevents independent lattice classes from carrying different invariant densities, and the absolute moment condition makes the displayed mean a well-defined vector.
\subsection{Short English statement}
Are all stationary laws of every irreducible mean-zero exclusion process mixtures of independent occupations with a common density?
\subsection{Sources}
\begin{itemize}
\item[S1] ULAM.AI and the UnsolvedMath Contributors, supplied problems.json, record 9600003 (AMR-095-0003), AMR Open Problem Lists. Catalogue identity and source transcription; CC BY 4.0.
\url{https://huggingface.co/datasets/ulamai/UnsolvedMath}.
\item[S2] Liggett - Some Open Problems (2012), Problem 3, PDF page 2. Original source indexed by AMR-095-0003.
\url{https://web.archive.org/web/20150919235903id_/http://www.math.ucla.edu/~tml/open.pdf}.
\end{itemize}
\end{document}
